\documentclass{article}
\usepackage{graphicx}
\usepackage{amsfonts}
\usepackage{amsmath}
\usepackage{amsthm}
\usepackage{tikz}

\theoremstyle{definition}
\newtheorem{theorem}{Theorem}[section]
\newtheorem{corollary}[theorem]{Corollary}
\newtheorem{proposition}[theorem]{Proposition}
\newtheorem{lemma}[theorem]{Lemma}
\newtheorem{definition}[theorem]{Definition}

\theoremstyle{remark}
\newtheorem{remark}[theorem]{Remark}
\newtheorem{example}[theorem]{Example}
\newtheorem{claim}[theorem]{Claim}

\title{Integral and the Fundamental Theorem of Calculus}
\author{Sarah Liu}
\date{January 2026}

\begin{document}

\maketitle

\section{Basics}
\subsection{Asymptotic Equivalence}
\[
x+1\sim x
\]
\[
\sinh x \sim \frac{1}{2}e^x
\]
\[
n!\sim e^{-n}n^n\sqrt{2\pi n}
\]
\begin{theorem}[Prime Number Theorem]
    Denote $\pi(x)$ as the number of prime number less than or equal to $\pi$.
    \[
    \pi(x)\sim\frac{x}{\ln x}
    \]
\end{theorem}


\end{document}